% =============================================================================
%  MONETARY BACKSTOP -- REVIEW COPY WITH TRACKED CHANGES (2026-10-07)
%
%  Red strike-through = deleted, yellow = added. Add to the preamble:
%
%    \usepackage{xcolor,soul,framed,cancel}
%    \colorlet{addcolor}{yellow!60}
%    \sethlcolor{addcolor}
%    \colorlet{shadecolor}{addcolor}
%    \renewcommand{\CancelColor}{\color{red}}
%    \newcommand{\yadd}[1]{\hl{#1}}                     % added text
%    \newcommand{\rdel}[1]{{\color{red}\st{#1}}}        % deleted text
%    \newcommand{\madd}[1]{\colorbox{addcolor}{$\displaystyle #1$}}  % added math
%    \newcommand{\mdel}[1]{{\color{red}\cancel{#1}}}    % deleted math
%    \newenvironment{yaddblock}{\begin{shaded}}{\end{shaded}}  % added paragraphs/equations
%    \soulregister{\eqref}{7}\soulregister{\cref}{7}
%    \soulregister{\cite}{7}\soulregister{\ref}{7}
%
%  Silent fixes (could not be shown as markup): $\m^{cb}$ -> $m^{cb}_t$ (\m is an
%  undefined control sequence); the Unicode minus in "$t − 1$" -> "$t-1$".
%
%  Labels eq:cb-objective, eq:cb-pl, eq:tpi-floor, eq:tpi-kkt are also defined in
%  02c-central-bank-tpi.tex: keep only one of the two sections in the paper.
%  Cites eq:hm-yield, eq:spread-def, eq:spread-prices, eq:yield-decomposition-D,
%  eq:spread-decomposition (02d) and eq:bond-pricing, eq:bond-decomposition,
%  eq:expectation (02-model).
% =============================================================================

\section{Monetary Backstop - WORK IN PROGRESS}
\label{sec:policy}

The central bank of the monetary union (CB) inhabits the \rdel{previously described environment}\yadd{environment described above} and conducts its unconventional monetary policy. The \rdel{central}\yadd{main} policy studied in this paper is the monetary backstop, \rdel{where}\yadd{under which} the CB assumes the role of \rdel{the}\yadd{a} buyer of last resort and intervenes in the bond market\rdel{s} to stabilise the macroeconomic situation.

\subsection{Monetary Backstop}
We \rdel{documented}\yadd{have shown} that an increase in sovereign risk transmits through the net worth of the intermediary, is recessionary\yadd{,} and endogenously raises risk premi\rdel{ums}\yadd{a}. The \rdel{beginning of the chain of reaction is, however simple, an}\yadd{chain of reactions, however, begins simply: an} increase in default risk lowers the price of the bond of the $D$ economy. To break the first \rdel{chain}\yadd{link} in that \rdel{long process}\yadd{chain}, the CB takes on a mandate to prevent \rdel{the rise of yield spreads}\yadd{the yield spread from rising above a cap}.

\paragraph{Main objective of the Central Bank:} In the baseline specification of the monetary backstop, \yadd{the} CB solves the following optimisation problem:
\begin{equation}
\min_{m^{cb}_t\,\ge\,0}\;m^{cb}_t
\qquad\text{s.t.}\qquad
\text{Spread}_t\le\;\overline{\text{Spread}}
\quad\text{if } d_t=0,
\qquad\qquad
m^{cb}_t=0\quad\text{if } d_t=1,
\label{eq:cb-objective}
\end{equation}
That is, the CB aims to minimise \rdel{the central bank}\yadd{its} purchases of the distressed bond, $m^{cb}_t$, subject to keeping the spread\yadd{, $\text{Spread}_t\equiv sp_t$ of \eqref{eq:spread-def},} below some threshold, $\overline{\text{Spread}}$, that \rdel{the central bank}\yadd{it} deems economically sound. The problem of the CB is \yadd{a} constrained optimisation\yadd{ problem}. \yadd{The} CB does not simply aim to minimise the yield spread; otherwise, it would purchase in every state and \rdel{would} turn the backstop \rdel{policy} into a permanent holder of the $D$ country's debt. The cap\yadd{,} as seen in \rdel{equation \ref{eq:cb-objective}}\yadd{\eqref{eq:cb-objective},} is what makes this policy\rdel{,} truly\rdel{,} a monetary backstop\yadd{: inactive inside the band and unlimited at its edge}.\\

\begin{yaddblock}
Because the flow yield \eqref{eq:hm-yield} is decreasing in the price, the cap is a
floor on the price of the $D$ bond. Writing $\bar c\equiv\overline{\text{Spread}}$,
\begin{equation}
q^{D}_t\;\ge\;\underline q_t\;\equiv\;\frac{\delta_b}{\delta_b+y\bigl(q^{F}_t\bigr)+\bar c}.
\label{eq:tpi-floor}
\end{equation}
Since purchases raise the market price of the bond, the solution of
\eqref{eq:cb-objective} is the complementarity condition
\begin{equation}
m^{cb}_t\;\ge\;0,
\qquad
q^{D}_t-\underline q_t\;\ge\;0,
\qquad
m^{cb}_t\,\bigl(q^{D}_t-\underline q_t\bigr)\;=\;0
\qquad\text{in } d_t=0 .
\label{eq:tpi-kkt}
\end{equation}
As long as the market clears the $D$ bond above the floor, the CB does nothing; when
the market would clear below it, the CB buys exactly the quantity that holds the
price at $\underline q_t$, and no more.
\end{yaddblock}

\rdel{There are f}\yadd{F}our \rdel{important} features of the objective \rdel{\ref{eq:cb-objective}}\yadd{\eqref{eq:cb-objective}} \rdel{that} matter for the analysis\rdel{ and are taken as assumptions}\yadd{; all four are assumptions}. First\rdel{ly}, the backstop \rdel{does not insure against default}\yadd{operates only in the no-default regime}. \rdel{CB will not buy a distressed asset that will for sure default. The backstop only insures against a priced default.}\yadd{The CB buys nothing once the sovereign has defaulted ($d_t=1$): it caps the price of default risk, not default itself. It does, however, bear the haircut on the bonds it already holds, so the banks that sold them to it are insured against that default.} Second\rdel{ly}, purchases are held to maturity\yadd{: $m^{cb}_t\ge0$, and bonds are never sold back}. This simplifies the problem for the central bank. Third, the floor moves with the safe yield. During flight-to-safety episodes\yadd{, $y(q^{F}_t)$ falls, $\underline q_t$ rises, and} the central bank must buy more to counteract this move. Fourth, the cap is exogenous and constant \rdel{in}\yadd{over} time\yadd{; in the baseline, $\bar c$ is 200 basis points a year}. \\

We \rdel{will} decompose the effects of the monetary backstop into two \rdel{effects}\yadd{parts by comparing two economies that are identical except for the backstop: in one, the purchase rule is part of the equilibrium and enters every expectation; in the other, it does not exist}. The announcement effect is the difference between the two economies at their stochastic rest points, where the backstop is in place but nothing is bought. The intervention effect is the difference between their responses to the same rise in default risk\yadd{, where the backstop} \rdel{where it} may be used. \rdel{So the}\yadd{Each response is measured from its own economy's rest point, so the} intervention \rdel{channel}\yadd{effect} works on top of the announcement effect.

\paragraph{\rdel{Sovereign Bond Sterilisation}\yadd{Sterilised purchases}:} The central bank, by assuming the role of a monetary backstop, takes on \rdel{the} risky assets in exchange for safe central bank paper. \rdel{The central bank}\yadd{It} pays for its bonds with a riskless one-period claim on itself, held by the $D$ banks and remunerated at the $D$ deposit rate\yadd{: interest-bearing reserves, in effect, much like the fixed-term deposits with which the ECB sterilised its Securities Markets Programme}. Let \rdel{the} $M^{cb}_t$ be the stock of \rdel{distressed} $D$ sovereign bonds \yadd{the CB} carr\rdel{ied}\yadd{ies} into period $t$. \rdel{Its value follows this law of motion:}\yadd{The book, the reserves outstanding at the end of $t$ and the obligation they create evolve as}
\begin{equation}
    M^{cb}_{t+1}=(1-\mdel{\delta_X^b}\madd{\delta_b})\,h_t\,M^{cb}_t+m^{cb}_t,
    \madd{\qquad
    \mathrm{res}_t=q^{D}_t\,M^{cb}_{t+1},
    \qquad
    O^{cb}_{t+1}=R_{D,t}\,\mathrm{res}_t ,}
\label{eq:cb-book}
\end{equation}
\yadd{so the book decays at the maturity rate and is written down by the same survival factor $h_t$ as every other holding: the CB ranks pari passu.}

\paragraph{Distribution of \rdel{seigniorage}\yadd{profits}:} The central bank remits the proceeds from its bond market operations\yadd{ in full every period, so it carries no net worth}. Its profit in period $t$ is the payoff on the book less the obligation on the reserves issued at $t-1$:
\begin{equation}
\Pi^{cb}_t=\left( \Xi^{D}_t- R_{D,t-1}q_{t-1}^{D}\right)M^{cb}_{t}
\madd{\;=\;\Xi^{D}_t\,M^{cb}_t-O^{cb}_t,
\qquad
\Xi^{D}_t=h_t\bigl[\delta_b+(1-\delta_b)\,q^{D}_t\bigr]}
\label{eq:cb-pl}
\end{equation}
This \rdel{seigniorage}\yadd{profit} is \rdel{then} distributed \rdel{back} to the governments of both countries\rdel{(}\yadd{ by the capital key: a share} $\kappa_D$ goes to the $D$ treasury and $1-\kappa_D$ to the $F$ treasury\yadd{. It is not seigniorage: the reserves pay the deposit rate, so the profit is the excess return of the bond over that rate. It is positive in expectation while the sovereign does not default, and negative when the sovereign defaults while the book is still held. With $\kappa_D=0.071$, 92.9\% of either goes to $F$.}

\rdel{The central bank, by assuming the role of a monetary backstop, takes on the risky assets in exchange for safe central bank paper.}


\subsection[Disorderly Markets and Uncertain Intervention]{\yadd{Disorderly Markets and Uncertain Intervention} - Still ideas not implemented}
\rdel{Disorderly Market, XYZ, and Uncertain Intermediaries}

We \rdel{will} analyse \rdel{three}\yadd{two} illustrative cases that highlight \rdel{the elements of the design of}\yadd{two design choices for} monetary backstops\yadd{: what the cap is measured against, and whether the market believes it}.
% the third case (XYZ in the old title) is left out until it is decided

\paragraph{Disorderly Market Interventions:} In this specification, \yadd{the} CB \rdel{operates under the assumption of risk neutrality of}\yadd{values the bond as if} financial intermediaries\yadd{ were risk neutral}. This yields an interpretation of excessive spreads as increases beyond \rdel{the expected payoff and liquidity constraint changes}\yadd{what the expected payoff and the tightness of the incentive constraint warrant}. The central bank \rdel{will} therefore aim\yadd{s} to correct only \rdel{for} the disorderly dynamics\yadd{, in the spirit of the ECB's TPI, whose stated aim is to counter unwarranted, disorderly market dynamics}.

\begin{yaddblock}
The CB's reference is the \emph{warranted price} of the $D$ bond: the value of its
payoff to a bank that faces the same deposit rate and the same constraint premium,
but is risk neutral across the default branches at every date,
\begin{equation}
q^{\star}_t\;=\;\frac{\E_t\bigl[h_{t+1}\bigl(\delta_b+(1-\delta_b)\,q^{\star}_{t+1}\bigr)\bigr]}
                   {R_{D,t}+\mathrm{lp}^{D}_t},
\qquad
\mathrm{lp}^{D}_t=\frac{\lambda_D\,\mu_{D,t}}{\E_t\bigl[\Omega_{D,t,t+1}\bigr]}.
\label{eq:q-warranted}
\end{equation}
This is \eqref{eq:bond-pricing} with the covariance term of
\eqref{eq:bond-decomposition} switched off at every date, today and in the
continuation. The cap is now set on the bond's excess over its own warranted yield,
not over the safe yield:
\begin{equation}
\min_{m^{cb}_t\,\ge\,0}\;m^{cb}_t
\qquad\text{s.t.}\qquad
y\bigl(q^{D}_t\bigr)-y\bigl(q^{\star}_t\bigr)\le\bar c^{\,dm}
\quad\text{if } d_t=0,
\qquad
m^{cb}_t=0\quad\text{if } d_t=1,
\label{eq:cb-objective-dm}
\end{equation}
so the floor of \eqref{eq:tpi-floor} becomes
$\underline q^{\,dm}_t\equiv\delta_b/\bigl(\delta_b+y(q^{\star}_t)+\bar c^{\,dm}\bigr)$
and \eqref{eq:tpi-kkt} holds with $\underline q^{\,dm}_t$ in place of $\underline q_t$.

What the CB now targets is the risk premium. Wherever the bank's pricing condition
holds with equality, the log form of \eqref{eq:spread-prices} applied to $q^{D}_t$
and $q^{\star}_t$ gives, to first order around a yield $y$,
\begin{equation}
y\bigl(q^{D}_t\bigr)-y\bigl(q^{\star}_t\bigr)
\;\approx\;\sum_{k\ge0}\,\frac{\delta_b+y}{1+y}\Bigl(\frac{1-\delta_b}{1+y}\Bigr)^{k}
\,\E_t\bigl[\mathrm{rp}^{D}_{t+k}\bigr],
\label{eq:dm-gap}
\end{equation}
a weighted average of the current and expected future risk premia of
\eqref{eq:yield-decomposition-D}, with weights that decay at the bond's maturity
rate and sum to one.

Three properties distinguish this rule from the baseline. First, a rise in
$\pi^{d}_t$ that only raises the expected loss $\mathrm{el}_t$, or a tightening of
the constraint that only raises $\mathrm{lp}^{D}_t$, lowers $q^{\star}_t$ and the
floor together with the market price, and triggers no purchase. The CB buys only
against the premium that arises because bankers value wealth most in the default
branch. Second, the floor no longer moves with the safe yield, so a flight to
safety does not by itself oblige the CB to buy. Third, the rule treats the
liquidity premium as warranted, although $\mathrm{lp}^{D}_t$ is the pass-through of
\cite{bocola2016pass}: by \eqref{eq:spread-decomposition}, mark-to-market losses
that tighten the constraint raise the sovereign spread and the credit spread one for
one. A disorderly-market backstop therefore lets this part of the balance-sheet
channel run. A variant that also counts the liquidity premium as disorderly drops
$\mathrm{lp}^{D}_t$ from the denominator of \eqref{eq:q-warranted}, which anchors
the cap to the actuarial value of the bond.

Unlike $\underline q_t$, which depends on the $F$ price alone, the floor
$\underline q^{\,dm}_t$ depends on the $D$ bank's own valuation, through
$\mu_{D,t}$ and $\E_t[\Omega_D]$, and on tomorrow's warranted price. The warranted
price is an equilibrium object that the rule itself moves.
\end{yaddblock}
% computation: q* is one more stored rule, q*(d,S), with the residual
% q*(R + lp) - E[h'(delta_b + (1-delta_b) q*')] = 0; every term is already formed in
% point_map._bond_demands (E_Om_D, mu_D, the regime weights and payD).

\paragraph{Uncertain Intervention: }Financial intermediaries price the backstop through their assessment of how likely it is to be used\yadd{. That probability has two parts: how likely the backstop is to be needed, which is endogenous, and how likely it is to be honoured when needed}. In the baseline, \rdel{assume that} the central bank is fully credible \rdel{in their policy}\yadd{: it honours the cap with certainty}. Here, we relax this assumption by introducing a \rdel{specific} credibility parameter\yadd{, $\vartheta\in[0,1]$,} that the intermediaries price \rdel{onto the agents}\yadd{into the bond}.

\begin{yaddblock}
Let $a_t\in\{0,1\}$ record whether the backstop is honoured at $t$. It is drawn
independently each period with $\Pr(a_t=1)=\vartheta$ and observed at the start of
the period. The CB solves \eqref{eq:cb-objective} when $a_t=1$ and stands aside
otherwise,
\begin{equation}
m^{cb}_t=
\begin{cases}
\text{the solution of \eqref{eq:cb-objective}}, & a_t=1,\\
0, & a_t=0,
\end{cases}
\label{eq:cb-uncertain}
\end{equation}
and the expectation operator \eqref{eq:expectation} gains a branch:
\begin{equation}
\E_t\bigl[g\bigr]
=\bigl(1-\pi^{d}_t\bigr)\,\E^{s}_t\Bigl[\vartheta\,g(0,1,s_{t+1})+(1-\vartheta)\,g(0,0,s_{t+1})\Bigr]
+\pi^{d}_t\,\E^{s}_t\bigl[g(1,s_{t+1})\bigr],
\label{eq:expectation-uncertain}
\end{equation}
where $g(d',a',\cdot)$ evaluates the decision rules of regime $(d',a')$. In default
nothing is bought, so $a'$ does not matter there. Let $\tilde q_{t+1}$ be the price
at which the market would clear at $t+1$ with no purchase at that date. Then
$q^{D}_{t+1}=\max\{\tilde q_{t+1},\underline q_{t+1}\}$ if $a_{t+1}=1$ and
$q^{D}_{t+1}=\tilde q_{t+1}$ if $a_{t+1}=0$, so the expected payoff in the
no-default branch is
\begin{equation}
\E^{s}_t\bigl[\Xi^{D}_{t+1}\,\big|\,d_{t+1}=0\bigr]
=\delta_b+(1-\delta_b)\,\E^{s}_t\bigl[\tilde q_{t+1}\bigr]
+\underbrace{\vartheta\,(1-\delta_b)\,\E^{s}_t\Bigl[\bigl(\underline q_{t+1}-\tilde q_{t+1}\bigr)^{+}\Bigr]}_{\text{the backstop's put, scaled by credibility}} .
\label{eq:put-uncertain}
\end{equation}
The last term is the value of a put option on the $D$ bond, written by the CB and
struck at the floor. The case nests both economies of the previous section:
$\vartheta=1$ is the baseline and $\vartheta=0$ the economy without a backstop.

Three consequences follow. First, the announcement effect rises with credibility.
The put enters today's price at every state, including states where the floor is far
from binding, and it enters with weight $\vartheta$. Second, a less credible
backstop is more expensive when it is used. With less of the put in the price,
$\tilde q$ is lower, the gap to the floor is wider, and the CB must buy more each
time the floor binds. Third, credibility risk is itself priced. Within the
no-default branch the bond now pays less when the backstop fails, $a_{t+1}=0$, in
exactly the states where the floor would bind and bankers value wealth most. The
covariance in \eqref{eq:bond-decomposition} gains a negative term and
$\mathrm{rp}^{D}_t$ rises: an uncertain backstop adds a premium of its own. The
probability that the CB intervenes next period becomes
\begin{equation}
\Pr_t\bigl(m^{cb}_{t+1}>0\bigr)
=\bigl(1-\pi^{d}_t\bigr)\,\vartheta\,
\E^{s}_t\Bigl[\mathbb{1}\bigl\{\mathcal{S}_{t+1}\in\mathcal{S}^{\rm bind}\bigr\}\Bigr],
\label{eq:prob-intervention-uncertain}
\end{equation}
where $\mathcal{S}^{\rm bind}$ is the set of states at which the floor binds.
A natural extension makes credibility state dependent, $\vartheta_t=\vartheta(s_t,M^{cb}_t)$,
decreasing in the risk factor and in the CB's exposure: a central bank that weighs
the expected loss on its book is least credible in exactly the states in which the
backstop is needed.
\end{yaddblock}
% computation: a_t is a regime, not a state. The rules are indexed by
% (d,a) in {(0,1),(0,0),(1,.)}: three regimes instead of two, so 1.5x the unknowns
% and ~2.25x the dense-Jacobian solve.
